\scanpage{211}
\mspage{127}
% archive sheet number: 210
% element: text | top, introduction | Ces conditions sont satisfaites
Ces conditions sont satisfaites dans les \unsure{trois|deux} cas suivants :
% element: text | upper third, A | pi power equality
\begin{enumerate}[label={},leftmargin=2em]
\item[A)] $\pi^{\pi_0}=\pi$ (pour mémoire, c'est contenu dans 2\textsuperscript{o})
% element: text | upper third, B | Les familles épimorphiques
\item[B)] Les familles épimorphiques strictes dans $M$ sont épimorphiques strictes universelles.
% element: margin note | upper left, alongside B | Il suffit de supposer
\margin{Il suffit de supposer que les \unsure{épimorph.} stricts soient \unsure{stricts universels|universels} ; cf \unsure{compléments} 4.5.18.}
% element: text | upper middle, C | Les limites inductives
\item[C)] Les limites inductives filtrantes dans \deleted{$\unsure{C|E}$}$M$ commutent au chgt de base (ce qui est le cas notamment si les $\varinjlim$ filtrantes sont exactes \deleted{\unsure{finies}} à gauche).
\end{enumerate}
% element: margin note | middle left, diagonal | Je soupçonne que
% slanted
\margin{Je soupçonne que cette condition \unsure{revient|équivaut} à la \unsure{$=$|même|$\simeq$} condition sur $C$, pour les $\varinjlim$ filtrantes de cardinal $\leq\pi$.}
% element: text | middle, Remarque | La filtration de M
\uline{Remarque} La filtration de $M$ par les $\mathrm{Filt}'{}^\pi$ n'est autre que celle envisagée dans SGA 4 I 9.13, où on fait de plus l'hypothèse B) sur $M$, qui après-coup n'était plus tellement de mon goût. Toujours est-il que d'après (4.5.13.1), cette filtration se rapproche plus de la ``bonne'' filtration par les $M_\pi$, que celle que j'avais \unsure{cru|eu} \unsure{avoir raison de} lui préférer, celle par les $\mathrm{Filt}^\pi$. Elle a \inserted{d'ailleurs} sur cette dernière l'avantage, \inserted{dans le cas envisagé B)}, d'être une filtration cardinale \inserted{(ce qui n'est pas \unsure{même le cas|vrai} pour $\mathrm{Filt}^\pi$,} \deleted{Mais je} \deleted{\unsure{sauf erreur}} du moins s'il existe des cardinaux infinis $\pi>\pi_0$ avec $\pi^{\pi_0}\unsure{\neq|\not\geq|\not\leq}\pi$). Je soupçonne \deleted{\unsure{que}} maintenant qu'en général (et s'il existe des cardinaux infinis $\pi>\pi_0$ tels que $\pi^{\pi_0}\neq\pi$), en dehors des conditions A) B) C), la filtration par les $(\mathrm{Filt}'{}^\pi)$ n'est \uline{pas} une
